%! Author = sriadityavedantam
%! Date = 3/25/26

\lesson{46}{Self-Study Two}{Day 46}

\begin{lemma}[Lebesgue Big Number Lemma]{
	If $(X,d)$ is sequentially compact and $\{U_i\}$ is an open cover of $X$, then there exists $r > 0$ such that for all $x\in X$, there exists $i$ with
	\[
		B_r(x)\subseteq U_i.
	\]
}
	Assume this is false.
	Then for every $r > 0$, there exists $x\in X$ such that for every $i$,
	\[
		B_r(x)\nsubseteq U_i.
	\]
	In particular, for each $n\in\n$, choose $x_n\in X$ such that
	\[
		B_{1/n}(x_n)\nsubseteq U_i
	\]
	for every $i$.
	Since $X$ is sequentially compact, there exists a subsequence $x_{n_k}\to x$.
	Since $\{U_i\}$ is an open cover, there exists $i_0$ such that $x\in U_{i_0}$.
	Because $U_{i_0}$ is open, there exists $r_0 > 0$ such that
	\[
		B_{r_0}(x)\subseteq U_{i_0}.
	\]
	Choose $K$ sufficiently large such that
	\[
		\frac{1}{n_K}<\frac{r_0}{2}
		\qquad\text{and}\qquad
		d(x,x_{n_K})<\frac{r_0}{2}.
	\]
	Now if $y\in B_{1/n_K}(x_{n_K})$, then
	\begin{align*}
		d(x,y)
		&\leq d(x,x_{n_K}) + d(x_{n_K},y)\\
		&< \frac{r_0}{2} + \frac{r_0}{2}\\
		&= r_0.
	\end{align*}
	So
	\[
		B_{1/n_K}(x_{n_K})\subseteq B_{r_0}(x)\subseteq U_{i_0},
	\]
	a contradiction.
\end{lemma}

\begin{definition}[Totally Bounded]
	A metric space $X$ is totally bounded if for all $\eps > 0$, there exist $y_1,\ldots,y_K\in X$ such that
	\[
		X\subseteq \bigcup_{i=1}^K B_\eps(y_i).
	\]
\end{definition}

\begin{lemma}{
	If a metric space is sequentially compact, then it is totally bounded.
}
	Assume $X$ is not totally bounded.
	Then there exists $\eps > 0$ such that no finite collection of $\eps$-balls covers $X$.
	Choose $x_1\in X$.
	Since $B_\eps(x_1)$ does not cover $X$, choose
	\[
		x_2\in X\setminus B_\eps(x_1).
	\]
	Inductively, having chosen $x_1,\ldots,x_n$, since
	\[
		\bigcup_{i=1}^n B_\eps(x_i)
	\]
	does not cover $X$, choose
	\[
		x_{n+1}\in X\setminus \bigcup_{i=1}^n B_\eps(x_i).
	\]
	Then for all $m\neq n$,
	\[
		d(x_m,x_n)\geq \eps.
	\]
	So $(x_n)$ has no convergent subsequence, contradicting sequential compactness.
	Hence $X$ is totally bounded.
\end{lemma}

\begin{definition}[Norm]
	A norm is a function
	\[
		\norm{\cdot}:V\to[0,\infty)
	\]
	such that for all $u,v\in V$ and $\lambda\in\rr$:
	\begin{enumerate}
		\item $\norm{v}\geq 0$, and $\norm{v}=0$ if and only if $v=0$,
		\item $\norm{\lambda v} = \abs{\lambda}\norm{v}$,
		\item $\norm{u+v}\leq \norm{u}+\norm{v}$.
	\end{enumerate}

\end{definition}

\begin{definition}[Support of $f$]
	The support of $f$ is
	\[
		\operatorname{supp}(f)\coloneqq \overline{\{x:f(x)\neq 0\}}.
	\]
\end{definition}

\begin{theorem}{
	A metric space is sequentially compact if and only if it is topologically compact.
}
	\textbf{($\implies$).}
	Suppose $X$ is sequentially compact.
	Let $\{U_i\}$ be an open cover of $X$.
	By the Lebesgue Big Number Lemma, there exists $r>0$ such that for every $x\in X$, there exists $i$ with
	\[
		B_r(x)\subseteq U_i.
	\]
	By total boundedness, there exist $y_1,\ldots,y_K\in X$ such that
	\[
		X\subseteq \bigcup_{j=1}^K B_r(y_j).
	\]
	For each $j$, choose $i_j$ such that
	\[
		B_r(y_j)\subseteq U_{i_j}.
	\]
	Then
	\[
		X\subseteq \bigcup_{j=1}^K U_{i_j}.
	\]
	So $X$ is compact.

	\textbf{($\impliedby$).}
	Suppose $X$ is compact.
	Let $(x_n)$ be a sequence in $X$.
	If the set $\{x_n:n\in\n\}$ is finite, then some value occurs infinitely often, giving a constant subsequence.
	So assume $\{x_n:n\in\n\}$ is infinite.
	If no subsequence converges, then for each $x\in X$ there exists $\eps_x>0$ such that
	\[
		B_{\eps_x}(x)
	\]
	contains at most finitely many terms of the sequence.
	Then
	\[
		\{B_{\eps_x}(x):x\in X\}
	\]
	is an open cover of $X$.
	By compactness, there is a finite subcover
	\[
		X\subseteq \bigcup_{j=1}^K B_{\eps_{x_j}}(x_j).
	\]
	Each ball contains only finitely many terms of the sequence, so altogether the sequence has only finitely many terms, a contradiction.
	Hence some subsequence converges.
	So $X$ is sequentially compact.
\end{theorem}

\begin{definition}[Lipschitz]
	A function $f:X\to Y$ is Lipschitz if there exists $K\geq 0$ such that for all $x,y\in X$,
	\[
		d_Y(f(x),f(y))\leq K\,d_X(x,y).
	\]
\end{definition}

\begin{proposition}{
	Lipschitz implies continuity.
}
	Let $\eps > 0$.
	If $K=0$, then $f$ is constant, hence continuous.
	If $K>0$, choose
	\[
		\delta \coloneqq \frac{\eps}{K}.
	\]
	Then
	\[
		d_X(x,y)<\delta \implies d_Y(f(x),f(y))\leq K\,d_X(x,y)<K\delta=\eps.
	\]
	So $f$ is continuous.
	In fact, $f$ is uniformly continuous.
\end{proposition}

\begin{definition}[Uniform Continuity]
	A function $f:X\to Y$ is uniformly continuous if for all $\eps > 0$, there exists $\delta>0$ such that
	\[
		d_X(x,y) < \delta\implies d_Y(f(x),f(y)) < \eps.
	\]
\end{definition}